\chapter*{Предисловие}
\addcontentsline{toc}{chapter}{Предисловие}

Настоящее учебное пособие предназначено для студентов курса прикладного статистического анализа данных на старших курсах бакалавриата и в магистратуре.
Материал издания основан на одноимённом лекционном курсе, читаемом на четвёртом курсе бакалавриата Московского физико-технического института (МФТИ).
Основная цель данного издания заключается не в систематизации формул и критериев, а в формировании навыка связного статистического рассуждения.
Читатель последовательно проходит путь от постановки задачи и визуального осмысления данных к обоснованному выбору метода, проверке предпосылок и содержательной интерпретации результата в предметной области.

На протяжении всего текста неоднократно обсуждаются три взаимосвязанных вопроса, без учёта которых численный вывод теряет научную состоятельность.
Первый касается того, что именно измеряется и какие допущения о случайности и независимости наблюдений принимаются в постановке.
Второй относится к применимости выбранной процедуры, к проверке её предпосылок и к тому, насколько выводы подкреплены графиками, а не одним лишь p-value.
Третий затрагивает величину и практическую значимость эффекта, а также правдоподобность альтернативных объяснений.
Такая структура изложения призвана приучить к единому стилю анализа данных, а не к набору разрозненных вычислительных рецептов.

Материал организован в виде последовательного повествования.
Каждая глава развивает одну тему от постановки к интерпретации, а иллюстративные примеры включены в текст рассуждения, а не сведены в отдельный перечень формул.

Изложение начинается с общего языка статистики, включая описание выборки, распределения, оценивание неопределённости и проверку гипотез в параметрической и непараметрической постановке.
Далее рассматриваются множественные сравнения и описание зависимостей между признаками, от корреляции и таблиц сопряжённости до сравнения групп в дисперсионном анализе.
Затем вводится регрессионное моделирование, включая линейные модели, их диагностику и обобщения для счётных, бинарных и иных типов отклика.
Следующий этап посвящён данным, упорядоченным во времени, разложению ряда, прогнозированию, стационарности и моделям авторегрессии со скользящим средним.
Завершающий теоретический раздел охватывает причинные графы и условную независимость, последовательные модели со скрытыми состояниями, байесовское обновление сведений о параметрах и сравнение моделей с учётом неопределённости.
Практические задания курса собраны в заключительной главе и воспроизводят указанную логику от первого знакомства с данными до самостоятельного отчёта, содержащего графики, обоснованный выбор метода и выводы, сформулированные на естественном языке.

Рекомендуется читать пособие в том порядке, в котором темы излагаются в семестровом курсе.
При повторном обращении к материалу удобно пользоваться таблицей обозначений и контрольными вопросами в конце теоретических глав.
В каждой теоретической главе раздел \emph{Практика и семинар} содержит разобранные задачи с иллюстрациями: постановка, графики, числа и вывод словами.
Повторить расчёты и сдать домашнее задание можно по материалам курса в репозитории PSAD (\PSADRepo{}); в книге язык программирования не фиксируется.

Текст шире слайдов лекций за счёт пояснений из лекционных материалов и разборов семинаров.

Пособие рассчитано на самостоятельную работу с текстом.
Табличные данные практических задач и их описание приведены в приложении.
Тексты, изображения и аудиозаписи, объём которых не позволяет включить их в книгу, указаны сносками к соответствующим задачам.
Для освоения теоретического материала и задач с табличными данными достаточно этого издания.
В конце каждой теоретической главы даны контрольные вопросы по теме, которые помогают закрепить определения, интерпретации и типичные ошибки постановки.
В заключительной части книги приведён список литературы по темам курса для тех, кто захочет углубить отдельные разделы.
